
  \item[(b)] \SGARefTarget{sga3:xvii:proposition:7:1:1:proof:ii:item:b}
    The group \({}_nE\) is finite over \(S\).  Indeed, there is an
    exact sequence
    \[
      0\longrightarrow{}_nH\longrightarrow{}_nE
      \longrightarrow{}_nK\xrightarrow{u}H_n,
    \]
    where \(H_n\) is the cokernel of the \(n\)-th power map on
    \(H\).  We know that \(H_n\) is of multiplicative type (\SGARefLink{sga3:ix:proposition:2:7}{Exposé
    IX, 2.7}), and hence separated, and that \({}_nH\) and
    \({}_nK\) are finite over \(S\).  The group \(\ker u\) is a
    closed subgroup of \({}_nK\), and hence is finite over \(S\);
    consequently \({}_nE\) is finite over \(S\), as an extension
    of one finite group by another (Exposé VI\textsubscript{B},
    \SGARefLink{prop:sga3-VIB-9.2}{9.2}).
  \end{enumerate}
\end{enumerate}

\SGASubsection{7.2}{Structure of Commutative Affine Algebraic Groups}
\SGARefTarget{sga3:xvii:subsection:7:2}

\SGAStatement{Theorem 7.2.1}{Let \(k\) be a field and \(G\) a
commutative affine algebraic \(k\)-group.  Then:
\begin{enumerate}
\item[(a)] \SGARefTarget{sga3:xvii:theorem:7:2:1:item:a}
  \(G\) contains a largest subgroup \(M\) of multiplicative type.
  The group \(M\) is characteristic in \(G\), the quotient \(G/M\)
  is unipotent, and the formation of \(M\) commutes with extension
  of the ground field \(k\).
\item[(b)] \SGARefTarget{sga3:xvii:theorem:7:2:1:item:b}
  If \(k\) is perfect, \(G\) is the direct product of \(M\) and a
  unipotent algebraic subgroup \(U\), uniquely so.
\end{enumerate}}
\SGARefTarget{sga3:xvii:theorem:7:2:1}

\SGAProof{Proof}
\begin{enumerate}
\item[(i)] Suppose that \(k\) is algebraically closed.  When \(G\)
  is smooth, \SGARefLink{sga3:xvii:theorem:7:2:1:item:b}{7.2.1(b)} is well known (BIBLE~4, Theorem~4).  If
  \(G\) is radicial of height \(1\), the decomposition of
  \(\operatorname{Lie}G\) described in \SGARefLink{sga3:xvii:lemma:4-2-2}{4.2.2} corresponds, in view
  of \SGARefLink{sga3:xvii:proposition:4:3:1:item:v}{4.3.1(v)} and Appendix II, \SGARefLink{sga3:ii:subsection:2-2}{2.2}, to a decomposition of \(G\) of
  the kind asserted in \SGARefLink{sga3:xvii:theorem:7:2:1:item:b}{7.2.1(b)}.  In general, \(G\) has a
  composition series whose successive quotients are smooth or
  radicial of height \(1\) (Appendix II, \SGARefLink{sga3:xvii:lemma:3-1}{3.1}).  To prove \SGARefLink{sga3:xvii:theorem:7:2:1:item:b}{7.2.1(b)},
  it therefore suffices to observe the following.  Suppose there
  is an exact sequence of commutative algebraic groups
  \[
    0\longrightarrow G'\longrightarrow G\longrightarrow
    G''\longrightarrow0,
  \]
  in which \(G'\), respectively \(G''\), is the product of a group
  of multiplicative type and a unipotent group,
  \[
    G'=M'\cdot U',\qquad G''=M''\cdot U''.
  \]
  Then the same is true of \(G\).  Indeed, consider the exact
  sequence
  \[
    0\longrightarrow G'/U'\longrightarrow G/U'
      \longrightarrow G''\longrightarrow0.
  \]
  By \SGARefLink{sga3:xvii:proposition:7:1:1:item:a}{7.1.1(a)}, the inverse image of \(M''\) in \(G/U'\) is a
  subgroup \(M_1\) of multiplicative type.  The group \(M_1\)
  lifts to a subgroup \(M\) of \(G\) by \SGARefLink{sga3:xvii:theorem:5:1:1:item:i}{5.1.1(i)}\SGARefLink{sga3:xvii:appendix:A}{(a)}.  Likewise,
  now using \SGARefLink{sga3:xvii:theorem:6:1:1:item:C}{6.1.1(C)}, one proves that there is a unipotent
  subgroup \(U\) of \(G\) which is an extension of \(U''\) by
  \(U'\).  Clearly \(G=M\cdot U\).

\item[(ii)] Let \(k\) now be arbitrary.  By \SGARefLink{sga3:xvii:subsection:6:4:item:i}{(i)},
  \(G_{\bar k}\) is the direct product of a group
  \(M_{\bar k}\) of multiplicative type and a unipotent group
  \(U_{\bar k}\).  Put
  \[
    S=\bar k\times_k\bar k,
  \]
  and let \(M_1\) and \(M_2\) be the two inverse images of
  \(M_{\bar k}\) under the two projections
  \(S\rightrightarrows\bar k\).  The group
  \[
    G_S/M_2=(G/M_{\bar k})_S
  \]
  has unipotent fibres; hence the image of \(M_1\) in \(G_S/M_2\)
  is trivial by \SGARefLink{sga3:xvii:proposition:2:4:item:i}{2.4(i)}, and \(M_1\subset M_2\).  The converse
  inclusion follows in the same way, so \(M_1=M_2\).  By fpqc
  descent, \(M_{\bar k}\) therefore descends to an algebraic
  subgroup \(M\) of \(G\).  It is clear that \(M\) is of
  multiplicative type, that \(G/M\) is unipotent, and that the
  formation of \(M\) is compatible with every extension of
  \(k\).  For every \(k\)-prescheme \(S\), every subgroup \(H\) of
  multiplicative type in \(G_S\) is contained in \(M_S\).
  Indeed, its image in the group with unipotent fibres
  \[
    (G/M)_S=G_S/M_S
  \]
  is trivial by \SGARefLink{sga3:xvii:lemma:2:5}{2.5}.
